\clearpage
\section{Checking for Functional Dissociation}
\label{app:dissociation}
\subsection{Goal heads and impulse heads}
The saved exploratory comparison uses 122 goal heads and 181 decision/impulse heads in Qwen2.5-Coder-32B. The sets share 36 heads (Jaccard overlap .135). Membership overlap alone does not establish whether the sets perform the same function.

\begin{figure}[ht]\centering
\includegraphics[width=.9\linewidth]{Figures/appendix/goal_vs_decision_overlap.png}
\caption{\textbf{Goal-head and impulse-head membership.} The saved head sets have partial overlap, concentrated in the later layers. Colors distinguish goal-only, impulse-only, and shared heads.}
\label{fig:goal-impulse-overlap}
\end{figure}

The accompanying intervention transfers a mean ``move-on'' signal into either set at the notebook's fixed pivot (relative position $-2$), using source coefficient 2. The evaluation uses 256 multiline prompts with five conditions and letter-named variables. One metric asks whether the full-vocabulary next-token argmax names the next condition; the other asks whether that token has higher probability than the current condition's token. These are next-token measurements, not GoalScope oracle scores or complete-program correctness.

\begin{figure}[ht]\centering
\includegraphics[width=\linewidth]{Figures/appendix/goal_vs_decision_moveon.png}
\caption{\textbf{Transferring a move-on signal through different head sets.} Existing results from \texttt{notebooks/clock/11\_goal\_vs\_impulse.ipynb}. Impulse-head intervention increases the move-on rate more than goal-head intervention in this setup. The sets differ in size and depth coverage, so this comparison does not isolate head identity at a matched budget.}
\label{fig:goal-impulse-moveon}
\end{figure}
\clearpage
